% !TeX root = ../教材题库_全部条目_转写工作稿.tex
% Source: XJ-B1-C06-TH-P215-Y1331
% 第6章 / 6.2 抽样 / 教材 p.215 / PDF 第222页
% OCR未校对：公式与文字均可能有误；未完成数学TeX转写；禁止直接发给学生
% 原题及相邻原解截图（核对用）：../题目原图/XJ-B1-C06-TH-P215-Y1331-01.webp
% 本文件只在转写完成后作为题干使用；OCR 条目尚不是可靠题干。
\begin{bankitem}{p.215 思考（含紧邻说明）}{第6章 / 6.2 抽样 / 教材 p.215 / PDF 第222页}{XJ-B1-C06-TH-P215-Y1331}{OCR未校对：公式与文字均可能有误；未完成数学TeX转写；禁止直接发给学生}
\begin{ocrdraft}
[教材p.215 / PDF第222页]
\par
？思考如果要了解育才中学高一年级1000名学生对科技、体育、艺术和劳技
\par
四类课外活动的喜爱情况，应如何进行抽样调查？
\par
为了使样本尽可能具有代表性，除抽取调查的学生数量要合适外，抽取样本时，
\par
不能偏向某些学生，应使该年级的每一名学生都有相等的机会被抽到．如果在抽样过
\par
程中，能使总体中的每个个体都有相同的可能性被选人样本，那么这样的抽样叫作随
\par
机抽样．随机抽样可以避免主观偏向，使样本具有代表性．
\par
人们经常用“任取”“随机抽取”或“等可能抽取”等来表示随机抽样．
\end{ocrdraft}
\end{bankitem}
